\documentclass[10pt,a4paper]{amsart}
\usepackage[margin=19mm]{geometry}
\usepackage{amsmath,amssymb,mathtools}
\usepackage[scr=rsfs,cal=euler]{mathalfa}
\usepackage{xfrac}
\usepackage[unicode,hidelinks,bookmarks=false]{hyperref}
\newcommand{\ie}{i.e.\ }
\newcommand{\cf}{cf.\ }
\newcommand{\resp}{respectively\ }
\newcommand{\Subsection}{\subsection}
\providecommand{\nobreakdash}{\nobreak\hbox{-}}
\setlength{\emergencystretch}{2em}
\multlinegap0pt
\usepackage{bm}
\usepackage{bbm}
\usepackage{dsfont}
\usepackage{xspace}
\usepackage{cancel}
\usepackage{mathtools}
\usepackage[shortlabels]{enumitem}
\usepackage{amsthm}
\makeatletter
\@ifundefined{theorem}{\newtheorem{theorem}{Theorem}}{}
\@ifundefined{lemma}{\newtheorem{lemma}[theorem]{Lemma}}{}
\makeatother
\makeatletter
\expandafter\ifx\csname subproof\endcsname\relax
\newenvironment{subproof}{\begin{trivlist}\item}{\end{trivlist}}
\fi
\def\to{\mathchoice{\longrightarrow}{\rightarrow}{\rightarrow}{\rightarrow}}
\makeatother
\title[Light bulb smoothing for topological surfaces in 4-manifolds]{Light bulb smoothing for topological surfaces in 4-manifolds}
\author{Jae Choon Cha, Byeorhi Kim}
\date{Source: arXiv:2303.12857v2 (CC BY 4.0)}
\begin{document}
\maketitle

\noindent\emph{Source and licence.} Light bulb smoothing for topological surfaces in 4-manifolds. Version arXiv:2303.12857v2 (CC BY 4.0), distributed under the Creative Commons Attribution 4.0 International licence, as recorded in the arXiv licence metadata for this exact version and shown on the abstract page https://arxiv.org/abs/2303.12857v2. Full rights notice: licenses/arxiv-2303.12857-CC-BY-4.0.txt.\par\smallskip

\noindent\textit{Editorial scope.} Counted unit: the Lemma whose source label is \texttt{lemma:pi\_1-null-approx} in the article (source0.tex lines 1046--1052, source label \texttt{lemma:pi\_1-null-approx}), delivered together with the article's own complete original proof of it (source0.tex lines 1061--1093, one \texttt{proof} environment of 33 lines). No mathematics is written, repaired or re-derived: statement and proof are the article's own text line for line. The proof is detached from the statement in the source: the article states the result at lines 1046--1052 and prints its proof at lines 1061--1093, and the proof environment's own header names the statement, so the association is the article's own. Required context retained uncounted: none. Every definition, display and label of those lines is retained unchanged. Omitted from this package and not redistributed: 1--1045, 1053--1060, 1094--2292. Nothing inside a retained excerpt is omitted or altered: every retained body is delivered exactly as the article's lines are written, so no omission receipt of any kind accompanies this package. Citations: the retained excerpts carry no citation command at all, so no bibliography entry of the article is transcribed and no reference is redistributed. Reference adaptation: none was needed and none was made; every reference inside the delivered unit resolves inside it, so no unresolved reference or citation is produced. Printed slips kept literal: no printed word, symbol, hypothesis or number of the counted statement or of its proof was altered. Layout and class: the article's own class is \texttt{amsart} with packages including geometry, float, array, multirow, calc, booktabs, hyperref, amssymb, amsxtra, bbm, enumitem, graphicx, color, tikz; the wrapper is the lane's pinned \texttt{amsart} editorial wrapper at 10pt on A4, which restates the theorem-like environments the retained text needs and adds the article's own macro definitions that the retained text needs (\texttt{\textbackslash to}), with extra wrapper packages (bm, bbm, dsfont, xspace, cancel, mathtools, enumitem). The delivered numbering is the unit's own. Assets: no retained span contains a figure environment, a graphics-inclusion command, a PDF-page inclusion, a table or a TikZ picture; no image file is copied into this package, the package declares no asset dependency and no image file is redistributed with it. Licence: version arXiv:2303.12857v2 of this article is distributed under the Creative Commons Attribution 4.0 International licence, as recorded in the arXiv licence metadata for this exact version and shown on the abstract page https://arxiv.org/abs/2303.12857v2. A full rights notice is delivered at licenses/arxiv-2303.12857-CC-BY-4.0.txt. Characters: every retained excerpt is pure ASCII, so the delivered engine, XeLaTeX with its default Latin Modern text font, reports no missing character.
\begin{lemma}
  \label{lemma:pi_1-null-approx}
  Let $K$ be a compact connected smooth manifold and $V$ be a topological manifold equipped with a metric~$d$.
  The manifolds $K$ and $V$ are allowed to have nonempty boundary.
  Let $f\colon K\hookrightarrow V$ be a topological embedding.
  Then there exists $\epsilon>0$ such that if a map $g\colon K\to V$ is $\epsilon$-closed to $f$, every double point of $g$ has a double point loop which is null-homotopic in~$V$.
\end{lemma}

\begin{proof}[Proof of Lemma~\ref{lemma:pi_1-null-approx}]
  Fix a metric $\rho\colon K\times K \to [0,\infty)$ induced by a Riemannian metric on~$K$.
  In $K$ and $V$, denote the metric ball with center $x$ and radius $r$ by $B(x,r)$.
  Use the compactness of $K$ to choose constants $r>0$ and $\delta>0$ such that 
  \begin{enumerate}[label=(\roman*)]
    \item\label{item:small-ball-contractible} $B(f(x),r)$ is contained in an contractible open set in~$V$ for all $x\in K$;
    \item\label{item:short-geodesic} there is a short geodesic between $x$ and $y$ in $K$ whenever $\rho(x,y)<\delta$.
  \end{enumerate}
  In addition, since $f$ is continuous, we may assume that
  \begin{enumerate}[resume*]
    \item\label{item:control-f} $d(f(x),f(y)) < r/2$ if $\rho(x,y)<\delta$.
  \end{enumerate}
  Use the continuity of $f^{-1}\colon f(K) \to K$ to choose $\delta' >0$ such that 
  \begin{enumerate}[resume*]
    \item\label{item:control-f-inverse} $\rho(x,y)<\delta$ if $d(f(x),f(y))<\delta'$.
  \end{enumerate}

  Let $\epsilon=\min(\delta', r)/2$.
  Suppose $g \colon K\to V$ is $\epsilon$-close to~$f$.
  If $g(x)=g(y)$, then
  \[
    d(f(x), f(y))\leq d(f(x),g(x))+d(g(y), f(y))<2\epsilon \le \delta'
  \]
  and hence $\rho(x,y)<\delta$ by~\ref{item:control-f-inverse}.
  There is a short geodesic $\alpha$ between $x$ and $y$ in~$K$ by~\ref{item:short-geodesic}.
  For any $z$ on $\alpha$, $\rho(x,z)\leq \rho(x,y)<\delta$, so $d(f(x),f(z)) < r/2$ by~\ref{item:control-f}.
  It follows that
  \[
    d(f(x),g(z)) \leq d(f(x),f(z))+d(f(z),g(z) )< r/2 +\epsilon \leq r.
  \]
  That is, $g(\alpha)\subset B(f(x), r)$.
  By~\ref{item:small-ball-contractible}, the double point loop $g(\alpha)$ is null-homotopic in~$V$.
\end{proof}

\end{document}
